%----------------------------------------------------------------------------------------
%	STAGE DESCRIPTION
%----------------------------------------------------------------------------------------
\section*{Scopo dello stage}
% Personalizzare inserendo lo scopo dello stage, cioè una breve descrizione
Lo scopo di questo progetto di stage è introdurre lo studente allo sviluppo di applicazioni web moderne in ambito SaaS, attraverso un percorso formativo guidato che parte dai fondamenti dello sviluppo full-stack e si estende fino ai concetti di architettura distribuita, containerizzazione e orchestrazione cloud adottati in \ragioneSocAzienda.

Lo studente avrà il compito di contribuire allo sviluppo di un modulo applicativo del gestionale interno aziendale, realizzando le funzionalità lato backend con NestJS 11 (su Node.js 24 LTS) e l'interfaccia lato frontend con Next.js 16 e Tailwind CSS 4, integrate a un database PostgreSQL 18. Nel corso del tirocinio, lo studente affiancherà il tutor aziendale nell'esplorazione di tematiche più avanzate quali l'architettura multitenant tipica delle piattaforme SaaS, il caching con Redis 8, la messaggistica asincrona con RabbitMQ 4 e l'utilizzo di un API Gateway (Kong Gateway 3.10 LTS). Verranno inoltre presentati gli ambienti di sviluppo (Docker Engine 29) e di produzione (AWS, Kubernetes 1.36 nelle distribuzioni k3s e k8s per i servizi backend, Vercel per il deploy dei frontend) già configurati dall'azienda, con l'obiettivo di fornire allo studente una visione completa del ciclo di vita di un'applicazione moderna, dallo sviluppo locale al deploy in cloud.

